\documentclass[11pt]{article}
\usepackage[utf8]{inputenc}
\usepackage[T2A]{fontenc}
\usepackage[russian]{babel}
\usepackage{amsmath,amssymb}
\usepackage[margin=2.5cm]{geometry}
\usepackage{hyperref}

\title{Методы агрегации между отправками}
\author{}
\date{11 сентября 2026}

\begin{document}
\maketitle

\section{Что агрегируется}

Стоит различать три независимых по сути объекта агрегации:
\begin{itemize}
  \item \textbf{Веса/дельты модели} — классический FedAvg
        (см.\ раздел \guillemotleft Основы FL\guillemotright).
  \item \textbf{Статистики} — например, квантили non-conformity score, а не веса
        модели; агрегируется скаляр или маленький вектор, не полная модель.
  \item \textbf{Эпизоды опыта} — сырые (или частично обработанные) траектории
        заливаются в общее хранилище, обучение идёт централизованно уже на них
        (fleet-scale RL, см.\ отдельный раздел про архитектуру флота).
\end{itemize}

\section{Как агрегируется}

\begin{description}
  \item[Простое (взвешенное) усреднение.] FedAvg; деградирует при сильном non-IID.
  \item[Secure aggregation.] Маскирование + Shamir Secret Sharing
        \cite{bonawitz2017} (см.\ отдельный раздел); сервер видит только сумму.
  \item[Робастная агрегация.] Median, trimmed mean, Krum и родственные схемы —
        защита не от подглядывания, а от byzantine-клиентов, шлющих отравленные
        обновления \cite{sable2023}.
  \item[Иерархическая агрегация.] Не все клиенты шлют напрямую на центральный сервер:
        сначала агрегация внутри кластера (edge-узел), затем кластеры агрегируются
        между собой \cite{asynchierarchical2022, ageofconvergence2023}. Снижает
        трафик и добавляет промежуточный уровень доверия/атаки.
  \item[Асинхронная агрегация.] Сервер не ждёт всех клиентов; обновление применяется
        сразу по прибытии, с весом, затухающим по staleness (насколько обновление
        \guillemotleft устарело\guillemotright) \cite{asyncsurvey2021}.
\end{description}

\section{Client--edge--cloud как конкретная схема}

Работа \emph{Timely Asynchronous Hierarchical Federated Learning: Age of
Convergence} \cite{ageofconvergence2023} формализует трёхуровневую схему: клиент
$\to$ edge-сервер (локальная агрегация внутри кластера) $\to$ центральный
cloud-сервер (глобальная агрегация). Для бытовых роботов это напрашивающаяся
топология: \guillemotleft дом $\to$ домовой edge-узел $\to$ облако вендора
\guillemotright, а не плоская звезда \guillemotleft все роботы $\to$ один
сервер\guillemotright.

\emph{Asynchronous Hierarchical Federated Learning} \cite{asynchierarchical2022}
уточняет механику: центральный сервер группирует воркеров в кластеры по топологии
сети, в каждом кластере выбирается локальный агрегатор; финальная агрегация
взвешивает вклады клиентов параметром staleness — решает проблему разноскоростных
устройств без блокировки на самого медленного.

Для геораспределённого флота (несколько регионов клиента) релевантна схема с
несколькими независимыми серверами вместо одного центрального
\cite{multiserver2024}.

\section{Применение к проекту}

Federated conformal prediction (секьюрная агрегация квантиля non-conformity score)
из исходной постановки проекта — частный случай \guillemotleft агрегации статистики,
а не весов\guillemotright{} из раздела выше. Стоит явно назвать это отдельным
агрегационным раундом с собственным бюджетом приватности, не смешанным с раундом
агрегации весов VLA-модели: у двух раундов разная чувствительность и разная частота.

\section{Открытый вопрос}

Меняет ли staleness-weighting профиль memorization-атаки на политику? Устаревшее
обновление может быть как менее, так и более информативным об исходных данных — в
зависимости от того, насколько модель успела \guillemotleft забыть\guillemotright{}
между раундами синхронизации. Это не закрыто ни в одном из найденных источников и
может стать самостоятельным пунктом (C4) поверх уже заявленных в security-роадмапе
C1--C3.

\section{Тезаурус}

\begin{description}
  \item[Взвешенное усреднение] Базовый способ агрегации в FedAvg: вклад каждого
    клиента усредняется пропорционально объёму его локальных данных.
  \item[Робастная агрегация] Семейство методов (median, trimmed mean, Krum),
    защищающих итоговый агрегат от небольшого числа испорченных (byzantine)
    вкладов, а не от подглядывания сервером.
  \item[Median / trimmed mean] Агрегация по медиане или по среднему после отбрасывания
    крайних значений — устойчивее к выбросам, чем простое среднее.
  \item[Krum] Алгоритм робастной агрегации, выбирающий на каждом раунде вклад,
    ближайший к большинству остальных, вместо усреднения всех.
  \item[Иерархическая агрегация] Схема, в которой клиенты сначала агрегируются внутри
    промежуточного кластера (например, edge-узла), и только затем кластеры
    агрегируются между собой на верхнем уровне.
  \item[Асинхронная агрегация] Сервер применяет обновление сразу по прибытии, не
    дожидаясь всех клиентов, с весом, ослабленным по staleness.
  \item[Staleness] Мера того, насколько обновление клиента \guillemotleft
    устарело\guillemotright{} к моменту применения — чем больше раундов прошло с
    момента его получения весов, тем выше staleness.
  \item[Client--edge--cloud] Трёхуровневая топология агрегации: клиент $\to$
    промежуточный edge-сервер $\to$ центральный облачный сервер.
  \item[Geo-distributed multi-server] Схема с несколькими независимыми серверами
    агрегации в разных географических регионах вместо одного центрального.
\end{description}

\begin{thebibliography}{9}
\bibitem{bonawitz2017} K.~Bonawitz et al. \emph{Practical Secure Aggregation for
  Privacy-Preserving Machine Learning}. ACM CCS, 2017.
\bibitem{sable2023} \emph{SABLE: Secure And Byzantine robust LEarning}.
  arXiv:2309.05395, 2023.
\bibitem{asynchierarchical2022} \emph{Asynchronous Hierarchical Federated Learning}.
  arXiv:2206.00054, 2022.
\bibitem{ageofconvergence2023} \emph{Timely Asynchronous Hierarchical Federated
  Learning: Age of Convergence}. arXiv:2306.12400, 2023.
\bibitem{asyncsurvey2021} \emph{Asynchronous Federated Learning on Heterogeneous
  Devices: A Survey}. arXiv:2109.04269, 2021.
\bibitem{multiserver2024} \emph{Asynchronous Multi-Server Federated Learning for
  Geo-Distributed Clients}. arXiv:2406.01439, 2024.
\end{thebibliography}

\end{document}
